\section{Persistance}

\subsection{Le trait \texttt{Repository}}

\texttt{pomone-db} expose un trait \texttt{Repository} qui couvre toutes les opérations dont les services \texttt{pomone-app} ont besoin : CRUD sur les \texttt{Crop}, \texttt{Variety}, \texttt{Planting}, \texttt{Location}, \texttt{Strata}, \texttt{Family}, \texttt{YearlyHarvest}, plus quelques requêtes agrégées.

Deux implémentations concrètes :

\begin{description}
  \item[SQLite] Cible mono-utilisateur, fichier local, idéal pour une exploitation individuelle (maraîchage, grandes cultures, arboriculture ou agroforesterie sur un poste isolé). Migrations dans \texttt{migrations/sqlite/}.
  \item[MariaDB] Cible multi-utilisateurs / serveur, pour les fermes collectives ou les déploiements partagés. Migrations dans \texttt{migrations/mariadb/}.
\end{description}

\subsection{Pourquoi deux backends}

Une partie du public visé n'a aucun serveur à disposition (SQLite suffit, base = fichier). Une autre partie veut partager les données entre plusieurs postes (MariaDB). Plutôt que d'imposer une seule contrainte, le trait \texttt{Repository} permet de servir les deux profils sans dupliquer la logique métier.

\subsection{Pas de logique dans les triggers}

Tentation classique en SQL : déclencher des calculs de dates ou de cycles dans des triggers. Pomone refuse cette approche pour trois raisons :

\begin{enumerate}
  \item Les triggers SQLite et MariaDB ont des comportements et des syntaxes différents : maintenir deux versions diverge vite.
  \item La logique métier en triggers est difficile à tester unitairement.
  \item La logique métier en Rust est typée, testable, et profite des newtypes du domaine.
\end{enumerate}

\subsection{Tests cross-backend}

Le fichier \texttt{cross\_backend\_tests.rs} paramètre les mêmes scénarios fonctionnels sur les deux implémentations. C'est ce qui garantit la parité de comportement.

\subsection{Migrations}

Les migrations sont versionnées (\texttt{0001\_initial.sql}, etc.) et appliquées au démarrage. Toute évolution de schéma passe par une nouvelle migration numérotée, jamais par une modification rétroactive d'une migration existante.

\subsection{Données de seed}

Les seeds (familles botaniques par défaut, strates agroforestières standards, jeu de données de démo) vivent dans \texttt{pomone-db/seed.rs}. Le seed de \emph{démo} a été retiré du runtime : il reste comme \emph{helper de test} pour reproduire un état réaliste dans les tests d'intégration.
